\section{Allowed lines and the eight open fan triangles}
\label{sec:area_law_cell_fan_sectors}

Fix an arbitrary origin, natural dyadic exponent and integer cell index.
Split every side of the actual square at its midpoint, giving eight
triangles $P_e$ with common center $c$. These are the elementary open
sectors used in the local construction of
\cite[Section~11]{OpenAI2026AreaLaw}. Put
$\pi_1(x)=x_1$, $\pi_2(x)=x_2$ and
$\Phi=\{\pi_1,\pi_2,\pi_1+\pi_2,\pi_1-\pi_2\}$.
The level lines $\phi(x)=\phi(c)$ are exactly the four allowed lines
through $c$.

% Source: 10-geometry.tex, prop:two-families, lines 299--310,
% and geometry:initial-stars, lines 352--363;
% revision adc7f1241b42e322a6451854ab7e4b4c146bf78a.
% This is the actual single-fan assertion, without a region assignment.

\begin{theorem}[Allowed lines miss open fan triangles]
    \label{thm:area_law_cell_fan_sector_directions}
    \lean{TNLean.PEPS.AreaLaw.Geometry.cellFanPolygon_interior_not_isAllowedSlope_sub_center}
    \leanok
    \uses{def:area_law_cell_fan,def:area_law_template}
    Every one of the eight actual triangles satisfies
    \begin{align}
        \operatorname{int}(P_e)\cap
        \bigcup_{\phi\in\Phi}\{x\in\mathbb R^2:\phi(x)=\phi(c)\}
         & =\varnothing.
    \end{align}
    Equivalently, the vector $x-c$ has none of the slopes
    $0,\infty,1,-1$ whenever $x\in\operatorname{int}(P_e)$.
\end{theorem}
\begin{proof}\leanok
    \uses{thm:area_law_unsplit_side_endpoints,
        thm:area_law_elementary_side_geometry}
    Let $\phi$ be any of the four surjective linear functionals.
    The whole-side coordinates and midpoint-half descriptions show
    that the two perimeter vertices $a_e,b_e$ lie in one weak half-plane
    through $c$: either both functional values are at least $\phi(c)$,
    or both are at most $\phi(c)$.

    In the first case, convexity places $P_e$ in
    $\phi^{-1}([\phi(c),\infty))$. Since a surjective continuous linear
    functional is open, the interior of this half-plane is
    $\phi^{-1}((\phi(c),\infty))$. Hence
    $\phi(x)>\phi(c)$ for every $x\in\operatorname{int}(P_e)$.
    Applying the same argument to $-\phi$ handles the other weak side.
    Thus none of the four level equations can hold in the interior.
\end{proof}
